%!TEX root = ../main.tex
% =============================================================================
%  第 9 章  问题 4：波动电价下的风险决策（9.1 / 9.2 / 9.3 三段式）
%  三段式固定：9.1 两层电价与联合情景→阈值律不变性→CVaR 前沿；9.2 求解 + 四日期结果 + 四问汇总；9.3 情景保真度与基线、三组对照与灵敏度、稳健性小结
%  本节含：① 电价过程的统计检验结论（完整统计量见附录 E）  ② 两层电价模型与联合情景
%          ③ 波动电价下的两阶段 LP（双重不确定性）  ④ 三组对照实验的结论
%          ⑤ CVaR 线性化与风险—费用前沿  ⑥ 求解结果
%  口径：C-11（CVaR）在本节首次定义；图名 fig08_q4_price.pdf / fig09_cvar_frontier.pdf
% =============================================================================

\section{问题 4：波动电价下的风险决策}\label{chap:q4}

\subsection{模型建立}

问题 4 中电价与光伏出力同时随机（分时电价时段划分见文献~\cite{ref1093}，日前计划与实时偏差的结算
机制见文献~\cite{ref1217}）。把附件 4 的电价序列直接当作“另一组已知价格”代入模型，等于假设
\(0{:}00\) 即知全天实际电价，与题意矛盾，故先把它当作\textbf{随机过程}检验：检验在 2025 年全年
（\(52\,560\) 个观测）上进行，费用结果仍只在 334 天上报告。四项检验的完整统计量与判定见
\appref{sec:price-test-app}{D.3 电价过程检验}，其建模含义可概括为三条。① \textbf{骨架稳定但残差不可
忽略}：附件 1 剖面与附件 4 逐时段均值的最大偏差仅 \(5.15\times10^{-5}\)~元/kWh，而残差逐时段标准差
均值达 \(0.1238\)~元/kWh，故采用“骨架 \(+\) 残差”两层结构。② \textbf{残差非正态且强自相关}：
JB 统计量 \(656.70\) 拒绝正态，日内 \(\mathrm{ACF}(1)=0.8589\)，峰段残差标准差为谷段的 2.81 倍，
故按时刻对齐分块重采样，不按独立同分布抽样。\textbf{该检验的边界须写明}：\(n=52\,560\) 时 JB 检验
对任何微小偏离都会拒绝，\textbf{单靠“拒绝正态”没有信息量}，故本文不把它当作分布假设的唯一依据——
真正支撑“不按独立同分布抽样”的是同一行的另两条证据（日内 \(\mathrm{ACF}(1)=0.86\) 的强自相关与
峰谷残差标准差比 \(2.81\)），以及 \appref{sec:q4-gen}{D.2 情景生成对照} 的分块 Bootstrap 对照
（不做自相关修正时
\(1\%\) 分位相对差达 \(20.4\%\)）。③ \textbf{电价与光伏的池化相关 \(+0.028\) 是假象}：
它被“夜间光伏恒为 0”稀释，日聚合尺度上的真实结构是 \(-0.343\)；若据此独立抽样，就会生成“午间光伏
大发、同时电价高企“这类不可能情景并\textbf{低估尾部风险}（图~\ref{fig:q4-price}）。

\begin{figure}[htbp]
  \centering
  \includefig[width=0.58\linewidth]{fig08_q4_price.pdf}
  \caption{电价过程的骨架与残差结构（上子图：附件 4 电价的逐槽分位带 10\%--90\% 与 25\%--75\%、
           逐槽中位数与附件 1 分时骨架；下子图：日内残差自相关 \(\mathrm{ACF}(1\text{--}6)\)，
           并标注日级残差自相关与峰谷波动比；范围：2025-01-01 至 12-31 共 365 天 \(\times\) 144 槽；
           单位：电价 元/kWh，自相关为无量纲量；样本：\(n=52\,560\) 个电价观测）}
  \label{fig:q4-price}
\end{figure}

\subsubsection{两层电价模型与联合情景生成}\label{sec:q4-price}

用\textbf{两层结构}刻画电价：第一层是可预测的\textbf{分时骨架} \(\bar p_t\)（取附件 1 的电价剖面，
全年尺度上等于附件 4 的逐时段均值），第二层是不可预测的\textbf{随机残差} \(\varepsilon_{d,t}\)：

\begin{equation}
  p_{d,t}=\bar p_t+\varepsilon_{d,t},\qquad \mathbb{E}[\varepsilon_{d,t}]=0 ,
  \label{eq:q4-price}
\end{equation}

其中 \(d\) 为日期、\(t\) 为槽位，单位为元/kWh。由检验②，残差须\textbf{按时刻对齐}重采样，故以
\textbf{整日}为单位构造情景。设 \(r_{d,t}\) 为第 \(t\) 槽的净需求预报残差，\textbf{符号与
\S\ref{sec:q2-scen} 式~\eqref{eq:q2-resid} 完全一致}，即
\(r_{d,t}=(\widehat{PV}^{\text{fcst}}_{d,t}-PV^{\text{act}}_{d,t})\Delta t\)（\(0{:}00\) 预报口径的
\emph{实际}净需求与\emph{预报}净需求之差，等价地“预报高估光伏为正”，kWh；与式~\eqref{eq:q4-lp-gap}
中的 \(+r_{\omega,t}\) 配套使用）；\(p_d\) 为该日平均电价；把 334 天按 \((\sum_tr_{d,t},\,p_d)\) 各分
8 层形成 \(8\times8\) 网格，每格取\textbf{medoid}（层内欧氏距离和最小的历史日）为代表，权重取层内样本
占比。所得 64 个代表日（权重 \(0.0030\)--\(0.0299\)，合计 1）保真两个维度的边缘分布，并因电价剖面与
净需求残差来自\textbf{同一个历史日}而\textbf{天然保留联合结构}（保真度明细见 \appref{sec:scenfid-app}{B.2 联合情景保真度检验明细}）。



对照口径为\textbf{分块 Bootstrap}（6~h 块、按时刻对齐）生成的 64 个合成情景：它不依赖历史日的可获取
性，但本题样本量下边缘保真度明显更差（均值偏差 \(3.41\%\)、\(1\%\) 分位相对差 \(20.4\%\)），故
\textbf{只作对照}。保留 \(\rho(p,\text{净负荷})\) 是必要的：该相关意味着“贵的时候恰好也在多买电”，
若独立抽样，“高价 \(\times\) 高净需求”的联合尾部占比会被人为压低，风险度量随之失去意义；二维网格取
代表日后 \(\rho=0.2597\)，与实际 \(0.2620\) 相差 \(0.0023\)，尾部结构被保真。

设 \(\Omega\) 为联合情景集（权重 \(w_\omega\)），\(r_{\omega,t}\) 为情景净需求残差，\(p_{\omega,t}\) 为
情景电价，则重算问题 2 的两阶段 LP 为

\begin{subequations}\label{eq:q4-lp}
\begin{align}
  \min\quad
  & \sum_{t=0}^{143}\bar p^{\text{sc}}_tE_t+\sum_{\omega\in\Omega}w_\omega\sum_{t=0}^{143}5p_{\omega,t}g_{\omega,t},
  \qquad \bar p^{\text{sc}}_t=\sum_{\omega\in\Omega}w_\omega p_{\omega,t} \label{eq:q4-lp-obj}\\
  \text{s.t.}\quad
  & E_t-c_t+d_t\ \ge\ \bigl(L_t-\widehat{PV}^{(0)}_t\bigr)\Delta t, && \forall t \label{eq:q4-lp-plan}\\
  & g_{\omega,t}\ \ge\ \bigl(L_t-\widehat{PV}^{(0)}_t\bigr)\Delta t+r_{\omega,t}+c_t-d_t-E_t,\quad g_{\omega,t}\ge0, && \forall t,\omega \label{eq:q4-lp-gap}
\end{align}
\end{subequations}

外加式~\eqref{eq:q3-soc} 的储电量递推、界与终端条件。两个不确定性分别进入\textbf{目标系数}
（电价 \(p_{\omega,t}\)）与\textbf{约束右端项}（净需求残差 \(r_{\omega,t}\)），即问题 4 相对问题 2 的
\textbf{双重不确定性}。决策时的电价信息取情景概率加权均值 \(\bar p^{\text{sc}}_t\)：附件 4 未提供日前
电价预报，故首阶段电价系数用代表日实测电价的\textbf{层加权平均}，它与分时骨架的最大偏差
\(0.0462\)~元/kWh（占均价 \(6.02\%\)）；两者都\textbf{不含当日实际电价}，故不构成数据泄漏（假设 A10）。
\textbf{问题 4-3 的三次调整子问题同样使用 \(\bar p^{\text{sc}}_t\)}，不使用当日实际电价；
\textbf{结算一律按附件 4 的当日实际电价}。问题 4-3 在
式~\eqref{eq:q4-lp} 上再叠加三次滚动修正，调整与违约费用按 \calref{C-09} 的偏差量定价结算。

\subsubsection{阈值律的形式不变性}

命题~\ref{prop:threshold} 只依赖“目标对购电量线性”与“储电量递推”两条结构性质，\textbf{不依赖电价的
分布假设}，故电价由确定变随机只改变阈值（水价值 \(\nu_t\)）的\textbf{取值}，不改变函数形式：
\(\alpha_t\le\eta_c\nu_t\) 时充电、\(\alpha_t\ge\nu_t/\eta_d\) 时放电、其间为死区
（两侧之比 \(1/(\eta_c\eta_d)=1.2346\)）。两点边界：① 结论\textbf{在储电量界与功率界不活跃时精确
成立}，界活跃时须回到式~\eqref{eq:q1-kkt} 并计入界乘子；② “形式不变”不等于“阈值是常数”——波动电价下
\(\nu_t\) 仍须\textbf{逐时段由场景对偶确定}，随情景集与风险参数 \(\lambda\) 变化 \evid{C-Q4-01}。

\subsubsection{CVaR 与风险—费用有效前沿}\label{sec:q4-frontier}

记 \(C_\omega\) 为情景 \(\omega\) 的总购电费用、\(w_\omega\) 为情景权重（分层代表日的层样本占比，
\(\sum_\omega w_\omega=1\)，非等权）、\(\alpha=0.95\) 为置信水平，条件风险价值
（Rockafellar--Uryasev 形式~\cite{refrockafellar}）为

\begin{equation}
  \mathrm{CVaR}_\alpha(C)=\min_{\eta}\Bigl\{\eta+\frac{1}{1-\alpha}
  \sum_{\omega\in\Omega}w_\omega\max\bigl(0,\,C_\omega-\eta\bigr)\Bigr\}.
  \label{eq:q4-cvar}
\end{equation}

\textbf{权重口径必须与情景集一致}：本文的情景集是分层 medoid 集，权重 \(0.0030\)--\(0.0299\)
相差约 10 倍，故式~\eqref{eq:q4-cvar} 用 \(w_\omega\) 加权；若误用等权（\(1/|\Omega|\)），
风险度量会系统性高估轻权重情景的作用，且与目标函数中期望项的加权口径不一致。本文的目标函数与
\textbf{报告的风险指标}都用同一组 \(w_\omega\)（\S\ref{sec:q4-frontier} 的前沿、\(\lambda^*\) 的
选取与表~\ref{tab:q4-summary} 的回测同源）。

引入 \(z_\omega\ge C_\omega-\eta\)、\(z_\omega\ge0\) 后，式~\eqref{eq:q4-cvar} 的分段线性结构完全由
线性约束表达，\textbf{加入风险度量后模型仍是 LP}，不引入整数变量。把 CVaR 作为风险成本嵌入微网调度
的做法见文献~\cite{refchenhan}。以 \(\lambda\in\{0,0.1,\dots,1.0\}\)（11 档）为参数扫描得
\textbf{风险—费用有效前沿}（图~\ref{fig:frontier}）；两端极端解为 \(\lambda=0\) 的纯期望最优与
\(\lambda=1\) 的纯风险最优，后者以 \(0.99\%\) 的期望费用换取 \(3.43\%\) 的 \(\mathrm{CVaR}\) 下降
（等权口径下同一组解报 \(3.19\%\)，差别只来自风险统计量的加权方式，见上段）。

\textbf{推荐点的判据事先锁定}：在“情景内期望费用相对 \(\lambda=0\) 上升不超过 \(0.5\%\)”
（门限 \(42\,205.57\) 元/日）下使 \(\mathrm{CVaR}_{0.95}\) 最小。满足该门限的最大档位是
\(\lambda^{*}=0.3\)：期望费用 \(42\,168.43\) 元/日（\(+0.41\%\)）、\(\mathrm{CVaR}_{0.95}\)
\(49\,358.05\) 元/日（\(-2.99\%\)）、\textbf{真实轨迹回测 \(45\,778.03\) 元/日（仅 \(+0.18\%\)）}，
紧急购电量由 \(480.19\) 降到 \(324.60\)~kWh/日（\(-32.40\%\)）。\textbf{拐点}：\(\lambda\) 从 0.7 升到
0.9 时 \(\mathrm{CVaR}_{0.95}\) 只再降 \(8.49\) 元/日（\(0.02\%\)），期望费用却继续上升
\(34.27\) 元/日（\(0.08\%\)）；\(\lambda=1\) 与 0.9 基本持平（\(49\,133.66\) 对 \(49\,135.01\)
元/日），即 \textbf{\(\lambda\approx0.7\) 以后风险已基本降不动}。三个代表点的情景内指标与真实轨迹
回测为：\(\lambda=0\) 时 \(41\,995.59\)、\(50\,877.95\)、\(45\,694.27\)；\(\lambda^{*}=0.3\) 时
\(42\,168.43\)、\(49\,358.05\)、\(45\,778.03\)；\(\lambda=1\) 时 \(42\,411.56\)、\(49\,133.66\)、
\(45\,952.85\)（元/日）。可见情景内的排序在\textbf{实际发生}的 334 天轨迹上同样成立，但幅度小得多
（\(+0.57\%\) 对 \(+0.99\%\)）——这正是把情景内最优直接当决策依据会高估收益的原因。

\begin{figure}[htbp]
  \centering
  \includefig[width=0.60\linewidth]{fig09_cvar_frontier.pdf}
  \caption{风险—费用有效前沿与真实轨迹回测（左子图：情景期望费用、\(\mathrm{CVaR}_{0.95}\)
           与真实轨迹回测费用随风险厌恶系数 \(\lambda\) 的变化；右子图：期望费用—\(\mathrm{CVaR}_{0.95}\)
           前沿散点，颜色表示 \(\lambda\)，星号标出边际风险下降最快的档位；
           对象：\(\lambda\in\{0,0.1,\dots,1.0\}\) 共 11 档的解；范围：2025-02-01 至 12-31 中每 6 天取 1 天的 56 个代表日；
           单位：费用 元/日；样本：64 个代表情景、11 个 \(\lambda\) 档位，\(\alpha=0.95\)）}
  \label{fig:frontier}
\end{figure}

\subsection{模型求解与结果分析}

\subsubsection{求解结果}

在波动电价下重解问题 2 与问题 3（表~\ref{tab:q4-summary}）。两问的费用口径与固定电价下\textbf{完全
相同}（\calref{C-03}、\calref{C-09}），差别只在电价序列取附件 4 实际值、决策电价信息取情景加权均值，
故两组结果可直接比较。

\begin{table}[htbp]
  \centering
  \caption{波动电价与固定电价下的结果对照（同一实际轨迹：2025-02-01 至 12-31 共 334 天；
           单位：费用 元，电量 kWh；波动价基线 \(17\,026\,782.86\)~元、固定价基线
           \(16\,407\,319.63\)~元，均为 \calref{C-01} 口径）}
  \label{tab:q4-summary}
  \footnotesize
  \setlength{\tabcolsep}{4pt}
  \begin{tabularx}{\linewidth}{L{3.35cm} C{2.35cm} C{2.05cm} C{2.15cm} Y}
    \toprule
    模型 & 总费用/\unit{元} & 相对基线节省率 & 紧急购电量/\unit{kWh} & 与固定价同结构的差 \\
    \midrule
    问题 2（固定价，两阶段） & 14\,264\,792.13 & 13.06\% & 172\,015.82 & --- \\
    问题 4-2（波动价，两阶段） & 15\,094\,255.84 & 11.35\% & 182\,759.86 & \(+829\,463.70\)（\(+5.81\%\)） \\
    问题 3（固定价，四次滚动） & 13\,932\,402.95 & 15.08\% & 64\,522.09 & --- \\
    问题 4-3（波动价，四次滚动） & 14\,722\,153.17 & 13.54\% & 63\,523.11 & \(+789\,750.22\)（\(+5.67\%\)） \\
    \bottomrule
  \end{tabularx}
\end{table}

\textbf{波动的成本。} 同一窗口与同一实际轨迹下，波动价下两阶段与四次滚动结构的费用分别高出
\(5.81\%\) 与 \(5.67\%\)，无储能基线高出 \(3.78\%\)，节省率因此由 \(13.06\%/15.08\%\) 降到
\(11.35\%/13.54\%\)。该上升\textbf{不能由价格水平解释}：窗口内附件 4 电价均值 \(0.7575\)~元/kWh
反而比方时骨架的 \(0.7662\)~元/kWh 低 \(1.13\%\)，基线费用仍高 \(3.78\%\)——差额来自\textbf{高电价与
高净负荷同时出现}（池化相关 \(+0.260\)），这正是“必须生成联合情景”的定量证据。

\textbf{滚动调整在波动电价下依然有效。} 四次预报把波动价下的费用降低 \(372\,102.66\) 元
（\(2.47\%\)），紧急购电量由 \(182\,759.86\) 降到 \(63\,523.11\)~kWh（\(-65.24\%\)）。
该降幅略大于固定价下的 \(2.33\%\)（\(0.14\) 个百分点），但\textbf{两口径的费用结构不同、样本只有一组}，
故只能表述为“本题参数下滚动调整在两种价格环境下都有效”，\textbf{不足以支撑“电价越不确定、滚动价值越大”
这一普遍结论}——要检验它需要按价格残差的方差分档重跑（本文未做）。
问题 4-3 的四项费用为：计划相关
\(13\,979\,692.49\)、违约 \(224\,406.57\)、超额 \(270\,998.49\)、紧急购电 \(247\,055.63\) 元，
合计 \(14\,722\,153.17\) 元（其中违约电量 \(664\,243.01\)~kWh、超额电量 \(259\,710.27\)~kWh）。
违约与超额费用同时高于固定价，说明\textbf{价格不确定使“调整到多少”本身也更难判断}。四次滚动下波动价
与固定价的紧急购电量几乎相同（少 \(998.98\)~kWh，\(-1.55\%\)），而两阶段结构下波动价多买
\(10\,744.04\)~kWh（\(+6.25\%\)），即\textbf{滚动修正几乎完全吸收了电价波动带来的额外缺口}。

\subsubsection{四个指定日期的结果与跨问对照}
四个指定日期的结果见下一小节（表~\ref{tab:key4}、表~\ref{tab:emg-q42}、表~\ref{tab:emg-q43}；
逐时段完整结果见 \texttt{result4-2.xlsx}、\texttt{result4-3.xlsx}）。

按题面要求，问题 4 只需给出相应计算结果（不限定表格格式），本文按与问题 2、3 相同的格式给出四个
指定日期的汇总与紧急购电明细，见表~\ref{tab:key4}、表~\ref{tab:emg-q42}、表~\ref{tab:emg-q43}；
四问主结果的横向对照见表~\ref{tab:summary}。

\begin{table}[htbp]
  \centering
  \caption{问题 4-2、4-3 在四个指定日期的汇总结果（单位：电量 kWh、费用 元；
           全部为当日 144 个槽的合计；“全天计划购电费”为 \calref{C-09} 偏差量定价口径下的
           \(C^{\text{plan}}\)（即计划相关项，\textbf{不含}违约与超额，与表~\ref{tab:q3-t1} 的
           “全天购电费”不是同一个量）；
           该日总费用 \(=\) 全天计划购电费 \(+\) 违约与超额相关费用 \(+\) 紧急购电费）}
  \label{tab:key4}
  \footnotesize
  \setlength{\tabcolsep}{3.5pt}
  \begin{tabular}{C{1.35cm} C{1.75cm} C{2.10cm} C{2.10cm} C{1.85cm} C{1.75cm} C{1.75cm}}
    \toprule
    方案 & 日期 & 全天计划购电量 & 全天调整购电量 & 全天计划购电费 & 紧急购电量 & 该日总费用 \\
    \midrule
问题 4-2 & 2025-03-20 & 74\,196.33 & --- & 46\,622.31 & 914.18 & 50\,923.43 \\
    问题 4-2 & 2025-06-21 & 35\,543.98 & --- & 17\,764.30 & 252.82 & 18\,587.09 \\
    问题 4-2 & 2025-09-23 & 74\,693.23 & --- & 49\,288.50 & 149.64 & 49\,839.31 \\
    问题 4-2 & 2025-12-21 & 99\,538.07 & --- & 76\,013.33 & 0.00 & 76\,013.33 \\
    \midrule
    问题 4-3 & 2025-03-20 & 74\,321.58 & 69\,229.87 & 43\,088.31 & 894.55 & 48\,781.27 \\
    问题 4-3 & 2025-06-21 & 35\,725.56 & 35\,440.19 & 17\,651.83 & 5.24 & 17\,769.08 \\
    问题 4-3 & 2025-09-23 & 74\,869.82 & 72\,305.05 & 47\,748.10 & 91.92 & 48\,876.69 \\
    问题 4-3 & 2025-12-21 & 99\,720.87 & 98\,953.92 & 75\,242.01 & 345.33 & 77\,242.42 \\
    \bottomrule
  \end{tabular}
\end{table}


\begin{table}[htbp]
  \centering
  \caption{问题 4-2 在四个指定日期的紧急购电时间段与购电量（题面表 3 格式；单位：kWh；
            时段按连续触发的 10~min 槽合并、左闭右开；取整口径同表~\ref{tab:emg-q2}；
            本方案在 \(2025\text{-}12\text{-}21\) 未触发紧急购电）}
  \label{tab:emg-q42}
  \footnotesize
  \setlength{\tabcolsep}{4pt}
  \begin{tabular}{C{1.85cm} L{9.30cm} C{1.85cm}}
    \toprule
    日期 & 紧急购电时间段（括号内购电量，kWh） & 合计（kWh） \\
    \midrule
    2025-03-20 & \(6{:}00\)--\(8{:}00\)（225.30）；\(8{:}10\)--\(8{:}40\)（55.07）；\(8{:}50\)--\(10{:}30\)（523.99）；\(10{:}40\)--\(11{:}20\)（109.82） & 914.18 \\
    2025-06-21 & \(4{:}00\)--\(4{:}50\)（5.24）；\(5{:}00\)--\(5{:}30\)（82.79）；\(14{:}00\)--\(14{:}30\)（164.79） & 252.82 \\
    2025-09-23 & \(6{:}20\)--\(6{:}30\)（1.30）；\(6{:}50\)--\(7{:}40\)（125.27）；\(11{:}00\)--\(11{:}10\)（23.07） & 149.64 \\
    2025-12-21 & 无 & 0.00 \\
    \bottomrule
  \end{tabular}
\end{table}

\begin{table}[htbp]
  \centering
  \caption{问题 4-3 在四个指定日期的紧急购电时间段与购电量（题面表 3 格式；单位：kWh；
            时段按连续触发的 10~min 槽合并、左闭右开；取整口径同表~\ref{tab:emg-q2}）}
  \label{tab:emg-q43}
  \footnotesize
  \setlength{\tabcolsep}{4pt}
  \begin{tabular}{C{1.85cm} L{9.30cm} C{1.85cm}}
    \toprule
    日期 & 紧急购电时间段（括号内购电量，kWh） & 合计（kWh） \\
    \midrule
    2025-03-20 & \(6{:}00\)--\(7{:}00\)（125.36）；\(7{:}40\)--\(9{:}20\)（181.60）；\(9{:}30\)--\(10{:}00\)（87.91）；\(10{:}50\)--\(12{:}00\)（325.44）；\(12{:}30\)--\(12{:}50\)（86.88）；\(13{:}20\)--\(13{:}40\)（87.36） & 894.55 \\
    2025-06-21 & \(4{:}00\)--\(4{:}50\)（5.24） & 5.24 \\
    2025-09-23 & \(7{:}00\)--\(7{:}40\)（61.40）；\(11{:}00\)--\(11{:}10\)（30.51） & 91.92 \\
    2025-12-21 & \(11{:}00\)--\(12{:}00\)（345.33） & 345.33 \\
    \bottomrule
  \end{tabular}
\end{table}


\begin{table}[htbp]
  \centering
  \caption{四问主要结果汇总（费用单位：元，电量单位：kWh；问题 1 为典型日，
           其余为 2025-02-01 至 2025-12-31 共 334 天；问题 4 的电价为附件 4 波动价；
           表中带 a 标记者见表注）}
  \label{tab:summary}
  \scriptsize
  \setlength{\tabcolsep}{3pt}
  \begin{tabularx}{\linewidth}{L{2.35cm} C{1.60cm} C{1.95cm} C{1.95cm} C{1.95cm} Y}
    \toprule
    指标 & 问题 1 & 问题 2 & 问题 3 & 问题 4-2 & 问题 4-3 \\
    \midrule
    计划购电量 & 59\,482.70 & 22\,220\,333.20 & 22\,220\,333.20 & 22\,198\,236.57 & 22\,255\,478.92 \\
    调整购电量 & --- & --- & 21\,843\,259.36 & --- & 21\,850\,946.24 \\
    充电量\textsuperscript{a} & 20\,740.67 & 89\,108.12 & 86\,965.89 & 88\,470.42 & 86\,423.35 \\
    放电量\textsuperscript{a} & 16\,799.94 & 72\,177.58 & 70\,442.37 & 71\,661.04 & 70\,002.91 \\
    紧急购电量 \calref{C-05} & --- & 172\,015.82 & 64\,522.09 & 182\,759.86 & 63\,523.11 \\
    盈余处置量 \calref{C-06} & 0.0000 & 3\,106\,390.84 & 2\,638\,963.35 & 3\,100\,463.79 & 2\,645\,427.81 \\
    计划购电费 & 35\,126.95 & 13\,637\,552.90 & 13\,216\,203.48 & 14\,397\,542.37 & 13\,979\,692.49 \\
    违约相关费用 \calref{C-09} & --- & --- & 210\,674.71 & --- & 224\,406.57 \\
    超额相关费用 \calref{C-09} & --- & --- & 262\,628.79 & --- & 270\,998.49 \\
    紧急购电费 & --- & 627\,239.24 & 242\,895.98 & 696\,713.47 & 247\,055.63 \\
    \textbf{总费用} & \textbf{35\,126.95} & \textbf{14\,264\,792.13} & \textbf{13\,932\,402.95} & \textbf{15\,094\,255.84} & \textbf{14\,722\,153.17} \\
    无储能基线 \(C^{base}\)（\calref{C-01}） & 48\,052.05（V1） & 16\,407\,319.63（V2） & 16\,407\,319.63（V2） & 17\,026\,782.86（V3） & 17\,026\,782.86（V3） \\
    相对 \(C^{base}\) 节省率（\calref{C-04}） & 26.90\% & 13.06\% & 15.08\% & 11.35\% & 13.54\% \\
    \bottomrule
  \end{tabularx}
  \normalsize
\end{table}

表注：标记 a 的\textbf{两行}（充电量、放电量）为\textbf{题面指定四个日期}（2025-03-20、06-21、
09-23、12-21）的合计值，逐日值见表~\ref{tab:q2-t2} 与表~\ref{tab:key4}；问题 1 列为典型日全天值。
\textbf{其余各行（含计划购电量）均为 334 天全窗口合计}。计划购电量按结果文件「全天购电量」汇总列
（模型内值）报出：问题 2 与问题 3 的该量\textbf{逐位相同}（\(22\,220\,333.2033\)~kWh），若改按 144 个
逐槽 4 位小数取整列求和则为 \(22\,220\,333.1391\)~kWh，相差 \(0.06\)~kWh（\(2.9\times10^{-9}\)），
属写入取整而非模型差异（\calref{C-02}）。表内各项\textbf{独立四舍五入}至 2 位小数，故“各项相加”与
“总费用”可相差 \(0.01\) 元（问题 2、3、4-3 各一处）；总费用按全精度值取整，两者的一致性判据是
\textbf{全精度偏差 \(<0.005\) 元}（\calref{C-03}），而非印制值列和为 0。紧急购电量 \calref{C-05}
分别在问题 2、3、4-2 的 296、302、294 天触发；问题 2 的 \(172\,015.82\)~kWh 占该窗口计划购电量的
\(0.77\%\)，对应紧急购电费占总费用 \(4.40\%\)。固定价口径下问题 2 日均节省 \(6\,414.75\) 元
（问题 1 典型日 \(12\,925.10\) 元），差额主要来自预报误差的紧急购电费（日均 \(1\,877.96\) 元）；
问题 3 四次滚动后升至 \(7\,409.93\) 元。\evid{C-Q1-03}\evid{C-Q3-01}

问题 1 的日费用可分解为三块（同口径、单位 元/日，式~\eqref{eq:q1-decomp}）：无储能基线
\(48\,052.05\)、储能削峰填谷收益 \(12\,925.10\)、最优购电费 \(35\,126.95\)；该分解是
\textbf{逐项恒等式}而非两个独立最优值之和。334 天窗口上，问题 2 的计划购电费与紧急购电费分别占
\(95.60\%\) 与 \(4.40\%\)——储能的价值在压低计划购电费，紧急购电费是预报误差的代价，两者不可
相抵后只报净数；四次滚动把四项占比改为 \(94.86\%\)、\(1.51\%\)、\(1.89\%\)、\(1.74\%\)。
波动价口径下问题 4-2、4-3 相对基线分别节省 \(1\,932\,527.02\) 元与 \(2\,304\,629.69\) 元，
\textbf{两套基线不得混用}。

\subsection{检验与灵敏度分析}

\subsubsection{内部一致性、情景保真度与基线对比}
\textbf{情景保真度。} 生成情景与历史实际的三项对照为：均值相对差 \(0.082\%\)、关键分位最大相对差
\(6.99\%\)、电价—净需求相关 \(0.2597\) 对实际的 \(0.2620\)，即\textbf{联合相依结构被保留}；
若独立抽样（丢弃该相关），会低估“高价又缺电”的联合尾部（\S\ref{sec:q4-price}）。
\textbf{内部一致性。} 情景保真度与恒等式残差、四项费用分解的自检结果与 \appref{sec:verify-three-app}{D.1 独立审计报告}
同口径；真实轨迹回测（图~\ref{fig:frontier} 右子图）与情景期望费用同向变化。
\textbf{基线对比。} 波动价下无储能基线为 \(17\,026\,782.86\) 元（变体 V3），问题 4-2、4-3 的费用为
\(15\,094\,255.84\) 与 \(14\,722\,153.17\) 元，分别节省 \(1\,932\,527.02\) 元（\(11.35\%\)）与
\(2\,304\,629.69\) 元（\(13.54\%\)）；\textbf{两套基线（固定价 V2 与波动价 V3）不得混用}：
波动价基线的费用比固定价基线高 \(3.78\%\)，而窗口内附件 4 的均价反而低 \(1.13\%\)，差额来自
电价与净负荷的共动，而不是价格水平。

\subsubsection{三组对照实验与灵敏度分析}

三组对照回答同一问题：\textbf{主结果是否依赖于某项建模选择}。完整表格与判定阈值见
\appref{sec:ctl-app}{D.2 对照实验明细}。

\textbf{① 情景生成方式不改变费用结论。} 情景集只用 2025-02-01 至 06-30（150 天）拟合，在 07-01 至
12-31（184 天）上按真实轨迹结算，四种生成器（历史代表日、分块 Bootstrap、KDE \(+\) Copula、
按光伏规模缩放）的样本外总费用极差 \(33\,355.98\) 元，仅为四者均值的 \(0.37\%\)，节省率落在
\(11.10\%\)--\(11.43\%\)。但\textbf{分项指标并不等价}：紧急购电量相差 \(48.5\%\)
（\(72\,324\)--\(107\,459\)~kWh），故引用该分项必须同时给出情景口径。据此主模型保留历史代表日生成器，
Copula 用于相依结构诊断与敏感性对照。\evid{C-Q4-03}

\textbf{② 风险是结构性的，须按目标整点配置。} 费用尾部集中在 \(15\)--\(18\) 时，缺电频率峰值出现在
\(11\) 时（\(23.8\%\)）与 \(5\) 时（波动价下 \(34.6\%\)），与 \S\ref{sec:data-err} 的预报偏差剖面一致，
故安全边际不应全天统一加成。

\textbf{③ 盒式分布鲁棒被随机规划全面占优。} 保守度 \(k\) 从 0 扫到 3.0，最优档 \(k^{*}=1.5\) 的回测
总费用 \(16\,819\,102.59\) 元，比两阶段随机规划高 \(11.43\%\)、\(\mathrm{CVaR}_{0.95}\) 高
\(8.65\%\)、紧急购电量高 \(79.6\%\)——\textbf{五项指标无一占优}，其价值只在于给出最坏情形上界。
\evid{C-Q4-04}

\textbf{关于已知负载的信息利用。} 正文 \S3.5 的 Copula 诊断显示电价与当日负载强正相依，而 R3
结构下当日负载已知；因此严格地说，首阶段电价系数应使用 \(E[p_t\mid L_d]\)，而不是无条件情景均值。
本文用历史代表日的日负载总量对情景重加权，以近邻与高斯核两种方式近似该条件期望，并在同一实际轨迹上
回测。结果显示三组条件化方案的费用反而比主模型高 \(0.05\%\)--\(0.44\%\)，说明在现有样本量与季节
异质性下，用历史负载代理条件期望会引入估计噪声；主模型保留 \(\bar p^{sc}_t\)，同时把条件化电价
预测列为部署校准项。完整对照见 \appref{tab:pricecond}{D.2 负载条件化电价预测对照}。\evid{C-Q4-06}

\subsubsection{稳健性小结}
① \textbf{费用结论对建模选择稳健}：四种情景生成器、负载条件化重加权、盒式鲁棒三组对照下，
“波动价使费用上升约 \(5.7\%\)--\(5.8\%\)、储能与滚动调整仍有效”的方向不变。
② \textbf{分项指标不稳健}：紧急购电量在四种生成器之间相差 \(48.5\%\)、在盒式鲁棒下高 \(79.6\%\)，
故该分项\textbf{必须与情景口径同时引用}。③ \textbf{一处结论外推的禁止}：滚动调整在波动价下多省
\(0.14\) 个百分点，但两口径费用结构不同、样本只有一组，\textbf{不能}据此得出“电价越不确定、滚动
价值越大”的普遍结论。④ \textbf{风险须按时段配置}：费用尾部集中在 \(15\)--\(18\) 时、缺电频率峰值在
\(11\) 时与 \(5\) 时，安全边际不宜全天统一加成。
